% Problem-section file following ../_template.tex.
\section{The PPT-squared conjecture}
\label{sec:problem-32}

\paragraph{Problem.}
Must the composition of any two compatible PPT completely positive maps be
entanglement breaking?  Let
$\Phi:M_{d_1}(\mathbb C)\to M_{d_2}(\mathbb C)$ and
$\Psi:M_{d_2}(\mathbb C)\to M_{d_3}(\mathbb C)$ be completely positive.  The
Choi operator of $\Phi$ is
\begin{equation}
  J(\Phi)
  :=\sum_{i,j=1}^{d_1}\lvert i\rangle\!\langle j\rvert
       \otimes\Phi(\lvert i\rangle\!\langle j\rvert),
  \label{eq:p32-choi-operator}
\end{equation}
and $J(\Psi)$ is defined analogously.  In the convention of
Eq.~\eqref{eq:p32-choi-operator}, a map $\Phi$ is PPT when
\begin{equation}
  (T\otimes\operatorname{id})\bigl(J(\Phi)\bigr)\succeq0,
  \label{eq:p32-ppt-condition}
\end{equation}
and it is entanglement breaking when $J(\Phi)$ is separable; the same
definitions apply to $\Psi$.  Under condition
Eq.~\eqref{eq:p32-ppt-condition} for both maps, is
$J(\Psi\circ\Phi)$ necessarily separable?  Equivalently, does postselection
on any joint measurement outcome on the middle systems of two PPT bipartite
states always leave a separable state on the two outer systems?

\subsection*{Status}

Unsolved

\subsection*{Source}

Christandl, M\"uller--Hermes, and Wolf formulate the PPT-squared statement as
an explicit conjecture about compositions of PPT maps
\sourcecite{ref:p32-christandl}{CMW19}.

\subsection*{Progress}

\begin{itemize}
  \item The conjecture was recorded in the Banff workshop report on operator
  structures in quantum information
  \sourcecite{ref:p32-birs}{BIRS12}.

  \item Christandl, M\"uller--Hermes, and Wolf formulated the two-map
  statement above, proved its equivalence to the self-composition version,
  and established it in equal dimension $d=2$, for Gaussian channels, and
  for several further special classes.  Their eventual
  entanglement-breaking results for repeated composition do not imply the
  two-composition claim \sourcecite{ref:p32-christandl}{CMW19}.

  \item The equal-dimension conjecture holds for $d=3$
  \sourcecite{ref:p32-chen}{CYT19}.  Consequently, equal dimension $d\geq4$
  is the first unresolved regime.

  \item The conjecture holds for diagonal-unitary-covariant maps, a class
  containing Choi-type, depolarizing, dephasing, and amplitude-damping
  examples \sourcecite{ref:p32-singh}{SN22}.  This covariance restriction
  does not cover arbitrary PPT maps.

  \item Other repeated-composition theorems prove eventual entanglement
  breaking for broad PPT families, but not after exactly two maps
  \sourcecite{ref:p32-kennedy}{KMP18}.

  \item A 2026 qutrit theorem proves a stronger composition statement for
  cones larger than the PPT cone and extends it to arbitrary selective
  measurements.  Its dimension-specific hypothesis leaves higher dimensions
  open \sourcecite{ref:p32-an-lee}{AL26}.
\end{itemize}

\subsection*{References}

\begin{enumerate}[label={},leftmargin=4.8em,labelsep=0.5em]
  \footnotesize
  \item[\textup{[BIRS12]}]\label{ref:p32-birs}
  M. B. Ruskai, M. Junge, D. Kribs, P. Hayden, and A. Winter (organizers),
  \emph{Operator Structures in Quantum Information Theory}, Banff
  International Research Station Workshop Report 12w5084 (2012).
  \href{https://www.birs.ca/workshops/2012/12w5084/report12w5084.pdf}{workshop report}.

  \item[\textup{[CMW19]}]\label{ref:p32-christandl}
  M. Christandl, A. M\"uller--Hermes, and M. M. Wolf, ``When Do Composed Maps
  Become Entanglement Breaking?'' \emph{Annales Henri Poincar\'e} \textbf{20},
  2295--2322 (2019).
  \href{https://doi.org/10.1007/s00023-019-00774-7}{doi:10.1007/s00023-019-00774-7};
  \href{https://arxiv.org/abs/1807.01266}{arXiv:1807.01266}.

  \item[\textup{[CYT19]}]\label{ref:p32-chen}
  L. Chen, Y. Yang, and W.-S. Tang, ``Positive-Partial-Transpose Square
  Conjecture for $n=3$,'' \emph{Physical Review A} \textbf{99}, 012337
  (2019). \href{https://doi.org/10.1103/PhysRevA.99.012337}{doi:10.1103/PhysRevA.99.012337};
  \href{https://arxiv.org/abs/1807.03636}{arXiv:1807.03636}.

  \item[\textup{[SN22]}]\label{ref:p32-singh}
  S. Singh and I. Nechita, ``The $\mathrm{PPT}^2$ Conjecture Holds for All
  Choi-Type Maps,'' \emph{Annales Henri Poincar\'e} \textbf{23}, 3311--3329
  (2022). \href{https://doi.org/10.1007/s00023-022-01166-0}{doi:10.1007/s00023-022-01166-0};
  \href{https://arxiv.org/abs/2011.03809}{arXiv:2011.03809}.

  \item[\textup{[KMP18]}]\label{ref:p32-kennedy}
  M. Kennedy, N. A. Manor, and V. I. Paulsen, ``Composition of PPT Maps,''
  \emph{Quantum Information and Computation} \textbf{18}, 472--480 (2018).
  \href{https://arxiv.org/abs/1710.08475}{arXiv:1710.08475}.

  \item[\textup{[AL26]}]\label{ref:p32-an-lee}
  J. An and S. Lee, ``Beyond the Positive Partial Transpose Squared
  Conjecture: The Qutrit Case,'' arXiv:2607.15947 (2026).
  \href{https://doi.org/10.48550/arXiv.2607.15947}{doi:10.48550/arXiv.2607.15947};
  \href{https://arxiv.org/abs/2607.15947}{arXiv:2607.15947}.
\end{enumerate}

\subsection*{Comment}

The two-composition claim is settled in equal dimensions two and three and
for several structured families.  No proof or counterexample is known for
arbitrary compatible finite-dimensional PPT maps, with equal dimension four
being the first open square case.

\subsection*{Tag}

Entanglement-breaking channels; Choi states; Partial transpose criterion;
Quantum channels

\subsection*{ID}

\texttt{op\_69520395226dc45a}
